\documentclass[a4paper]{article}
\usepackage[affil-it]{authblk}
\usepackage[english]{babel}
\usepackage[backend=bibtex,style=numeric]{biblatex}

\usepackage{geometry}
\geometry{margin=1.5cm, vmargin={0pt,1cm}}
\setlength{\topmargin}{-1cm}
\setlength{\paperheight}{29.7cm}
\setlength{\textheight}{25.3cm}

% useful packages.
\usepackage{amsfonts}
\usepackage{amsmath}
\usepackage{amssymb}
\usepackage{amsthm}     
\usepackage{enumerate} %enumerate environment
\usepackage{graphicx}
\usepackage{multicol}
\usepackage{fancyhdr}
\usepackage{layout}
\usepackage{ctex}      %中文环境
\usepackage{booktabs}  %三线表
\usepackage{verbatim}  %comment environment
\usepackage{float}     %强制表格图片位置
\usepackage{listings}  %insert code in latex
\usepackage{fontspec}  %font support
\usepackage{tikz}
\usepackage{makecell}
\usetikzlibrary{intersections}
\usepackage{arydshln}

\usepackage[colorlinks,linkcolor=blue]{hyperref}
\usepackage{xcolor}
\definecolor{mygreen}{rgb}{0,0.6,0}
\definecolor{lightgrey}{rgb}{0.95,0.95,0.95}
\definecolor{winered}{rgb}{0.5,0,0}
\definecolor{mymauve}{rgb}{0.58,0,0.82}
\lstset{
      backgroundcolor=\color{lightgrey},
      basicstyle             = \ttfamily,
      breakatwhitespace      = false,
      breaklines             = true,
      captionpos             = none,
      commentstyle           = \color{mygreen}\bfseries,
      extendedchars          = false,
      keepspaces             = true,
      keywordstyle           = \color{winered},         % keyword style
      language               = C++,                     % the language of code
      otherkeywords          = {string},
      numberstyle            = \tiny\color{mygray},
      rulecolor              = \color{black},
      showspaces             = false,
      showstringspaces       = false,
      showtabs               = false,
      stepnumber             = 1,
      stringstyle            = \color{mymauve},         % string literal style
      tabsize                = 2,
      title                  = \lstname,
      aboveskip              = 0pt,
      belowskip              = 0pt,
      columns                = flexible
}

% some common command
\newcommand{\dif}{\mathrm{d}}
\newcommand{\avg}[1]{\left\langle #1 \right\rangle}
\newcommand{\norm}[1]{\left\| #1 \right\|}
\newcommand{\difFrac}[2]{\frac{\dif #1}{\dif #2}}
\newcommand{\pdfFrac}[2]{\frac{\partial #1}{\partial #2}}
\newcommand{\OFL}{\mathrm{OFL}}
\newcommand{\UFL}{\mathrm{UFL}}
\newcommand{\fl}{\mathrm{fl}}
\newcommand{\op}{\odot}
\newcommand{\Eabs}{E_{\mathrm{abs}}}
\newcommand{\Erel}{E_{\mathrm{rel}}}
\addbibresource{citation.bib}

\begin{document}
% =================================================
\title{Numerical Analysis chapter 1 theoretical homework}

\author{Zhouyue Qian 3220104074
  \thanks{Electronic address: \texttt{3220104074@edu.zju.cn}}}
\affil{(Information and Computational Science Class 2201), Zhejiang University }


\date{Due time: \today}

\maketitle

%\begin{abstract}
%    The abstract is not necessary for the theoretical homework, 
%    but for the programming project, 
%    you are encouraged to write one.      
%\end{abstract}





% ============================================
\section*{I}
\subsection*{I-question}
Consider the bisection method starting with the initial interval $[1.5, 3.5]$. In the following questions “the
interval” refers to the bisection interval whose width
changes across different loops.
\begin{itemize}
    \item What is the width of the interval at the $n$th step?
    \item What is the supremum of the distance betweenthe root $r$ and the midpoint of the interval?
\end{itemize}
\subsection*{I-a}

%Give your answers here....\cite{Zhang2020BooleanAO}
If the initial width of the interval is $w$, then the width of the interval at the nth step will be $w_n = \frac{w}{2^n}=\frac{1}{2^{n-1}}$.

\subsection*{I-b}
This is determined by the initial width of the interval and is equal to $\frac{3.5-1.5}{2}=1$

at the very beginning, which is 1. This value will decrease as the bisection method narrows down the interval.


\section*{II}
\subsection*{II-question}
In using the bisection algorithm with its initial interval as $[a_0, b_0]$ with $a_0 > 0$, we want to determine the root
with its $relative$ error no greater than $\epsilon $. Prove that this goal of accuracy is guaranteed by the following choice 
of the number of steps,$$n \geq \frac{log(b_0-a_0)-log\epsilon -log(a_0)}{log2} -1$$

\subsection*{proof}

As we all know, the $relative$ error $\eta = \frac{|r - c|}{|r|} $,where $r$ is the true root, and $c$ is the approximate root found by the algorithm. 
The relative error being no greater than $\epsilon$ means that $\eta \leq \epsilon$, which equals to $|r - c| \leq \epsilon |r|$.
Also, we have already learnt that the width of the interval at the $n$th step is  $ \frac{b_0 - a_0}{2^{n+1}}$, where $n$ is the number of steps.
Therefore, we have $\frac{b_0 - a_0}{2^n}\leq|r - c| \leq \epsilon |r|  $Assuming $|r|=a_0$ (Since $a_0$ is a reasonable initial estimate), we have:
$$\frac{b_0 - a_0}{2^{n+1}}\leq|r - c| \leq \epsilon a_0$$
$$\frac{b_0 - a_0}{2^{n+1}}\leq \epsilon a_0$$
So we have: 
$$n \geq \frac{log(b_0-a_0)-log\epsilon -log(a_0)}{log2} -1$$ 
The proof is completed.  $\Box$ 

\section*{III}
\subsection*{III-question}
Perform four iterations of Newton's method for the
polynomial equation $p(x) = 4x^{3}-2x^{2}+3 = 0$ with
the starting point $x_0 = −1$. Use a hand calculator and
organize results of the iterations in a table.

\subsection*{III-answer}
The derivative of $p(x)$ is $p'(x) = 12x^{2}-4x$. The Newton's method iteration formula is:
$$x_{n+1} = x_n - \frac{p(x_n)}{p'(x_n)}$$
\begin{table}[ht]
  \centering
  \begin{tabular}{|c|c|}
  \hline
  Iteration & Value of $x_n$ \\ \hline
  $x_1$ & -0.8125    \\ \hline
  $x_2$ & -0.770804  \\ \hline
  $x_3$ & -0.768832  \\ \hline
  $x_4$ & -0.768828  \\ \hline
  \end{tabular}
  \caption{Iterations of $x_n$}
  \label{table:iterations}
  \end{table}
The answer was generated by $Python$ which helped by AI. \cite{kimiassistant}

\section*{IV}
\subsection*{IV-question}
Consider a variation of Newton's method in which only the derivative at $x_0$ is used,$$x_{n+1} = x_n - \frac{f(x_n)}{f'(x_0)}$$
Find $C$ and $s$ such that $$e_{n+1} = C(e_n)^s$$ where $e_n$ is the error of Newton's method at step n, s is
a constant, and C may depend on $x_n$, the true solution $\alpha$ , and the derivative of the function $f$.

\subsection*{IV-answer}
\begin{itemize}
  \item $Taylor \, Expansion$
    $$f(x) = f(x) + f'(\alpha )(x - \alpha ) + \frac{f''(\alpha )}{2!}(x - \alpha )^2 + \cdots$$
    Since \(f(\alpha) = 0\), we have:
    $$ f(x_n) = f'(\alpha)(x_n - \alpha) + \frac{f''(\alpha)}{2}(x_n - \alpha)^2 + \cdots $$
    $$ f(x_n) \approx f'(\alpha)(x_n - \alpha) + \frac{f''(\alpha)}{2}(x_n - \alpha)^2 $$
  \item $Error \; in \; Iteration \; Formula$
    $$ x_{n+1} = x_n - \frac{f(x_n)}{f'(x_0)} $$
    $$ x_{n+1} = x_n - \frac{f'(\alpha)(x_n - \alpha) + \frac{f''(\alpha)}{2}(x_n - \alpha)^2}{f'(x_0)} $$
    $$ x_{n+1} - \alpha = x_n - \alpha - \frac{f'(\alpha)(x_n - \alpha)}{f'(x_0)} $$
    $$ e_{n+1} = e_n - \frac{f'(\alpha) e_n}{f'(x_0)} $$
  \item $calculate$
    $$ e_{n+1} = (1 - \frac{f'(\alpha)}{f'(x_0)} ) e_n $$
    So, we could know that $C = 1 - \frac{f'(\alpha)}{f'(x_0)} $ and $s = 1$
\end{itemize}

\section*{V}
\subsection*{V-question}
Within $(-\frac{\pi }{2},\frac{\pi }{2})$, will the iteration $x_{n+1} = tan^{-1} x_n$ converge?

\subsection*{V-answer}
First, we note that the function \(\arctan(x)\) is the inverse tangent function, and its derivative is given by:
\[
\frac{d}{dx} \arctan(x) = \frac{1}{1+x^2}
\]
This derivative is always less than 1 for all \(x\) in the interval \((-\frac{\pi}{2}, \frac{\pi}{2})\), i.e.,
\[
\left|\frac{1}{1+x^2}\right| < 1
\]
According to $Contraction \; Mapping \; Theorem$ we could learn that $x_{n+1} = tan^{-1} x_n$ is converge and the fixed point is $(0,0)$.

\section*{VI}
\subsection*{VI-question}
Let \( p > 1 \). What is the value of the following continued fraction?
\[
x =  \cfrac{1}{p + \cfrac{1}{p  + \cfrac{1}{p + \cdots}}}
\]
Prove that the sequence of values converges. (Hint: this can be interpreted as \( x = \lim_{n \to \infty} x_n \), where
\[
x_1 = \frac{1}{p}, \quad x_2 = \frac{1}{p + \frac{1}{p}}, \quad x_3 = \frac{1}{p + \frac{1}{p+\frac{1}{p}}}, \quad \ldots, \quad \text{and so forth.}
\]
Formulate \( x \) as a fixed point of some function.)
\subsection*{VI-answer}
This continued fraction can be written as: $f(x) = \frac{1}{p+x}$; $x = f(x)$, which is a fixed point of the function $f(x)$. 
It is obvious that $\forall i,{x_i}>0$
The derivative of $f(x)$ is:
$$f'(x) = -(p+x)^{-2}<0$$ Since $p>1,x_1 = \frac{1}{p}<1$,
so we have $x_{n+1}<x_n<1$, which make sure that $|f'(x)<\frac{1}{p^{2}}<1|$,
so the sequence of values converges.

\section*{VII}
\subsection*{VII-question}

What happens in $problem II$ if $a_0 < 0 < b_0$? Derive
an inequality of the number of steps similar to that in
$II$. In this case, is the relative error still an appropriate
measure?

\subsection*{VII-answer}
Same as $problem II$, we have: the $relative$ error $\eta = \frac{|r - c|}{|r|} $,where $r$ is the true root, and $c$ is the approximate root found by the algorithm. 
We assume that $\alpha $ is the true root,$\epsilon $ is error, and $x_n$ is the approximate root found by the algorithm.
Also from $problem I$, we have:$x_n-\alpha \leq  \frac{|b_0 - a_0|}{2^{n+1}}$, for $\eta < \epsilon $ so we have:
$$\frac{|b_0 - a_0|}{2^{n+1}*|\alpha |} \leq \epsilon $$
$$ n \geq \frac{log(b_0-a_0)-log\epsilon -log(|\alpha |)}{log2} -1$$
The relative error cannot be an appropriate measure, since the estimate of $n$ in this expression depends on $\alpha $ .
So,we should use the absolute error instead of the relative error.

\section*{VIII}
\subsection*{VIII-question}
$(*)$ Consider solving $f(x) = 0 (f ∈ C^{k+1})$ by Newton's
method with the starting point x0 close to a root of
multiplicity k. Note that $\alpha$  is a zero of multiplicity $k$
of the function $f$ if
$$f^{(k)}\neq 0; \forall i< k, f^{(i)}(\alpha ) = 0$$
\begin{itemize}
    \item How can a multiple zero be detected by examining
    the behavior of the points (xn, f(xn))?
    \item Prove that if r is a zero of multiplicity k of the
    function f, then quadratic convergence in Newton's iteration will be restored by making this
    modification:
    $$x_{n+1}=x_{n}-k\frac{f(x_n)}{f'(x_n)}$$
\end{itemize}

\subsection*{VIII-answer}
\subsubsection*{question1}
Because $\alpha $ is a zero of multiplicity $k$ of the function $f$, we assume that $f(x)= (x - \alpha )^{k}g(x)$ and $g(\alpha )\neq 0 $
From the Newton's method iteration formula, we have:
$$x_{n+1}= x_{n}- \frac{f(x_n)}{f'(x_{n})}$$
$$x_{n+1}-\alpha = x_{n}-\alpha - \frac{f(x_n)}{f'(x_{n})}$$
$$x_{n+1}-\alpha = (x_{n}-\alpha)-\frac{(x_{n}-\alpha)g(x_n)}{kg(x_n)+(x_n-\alpha)g'(x_n)}$$
Now we assume $e_{i}=x_{i}-\alpha$, and we have:
$$\frac{e_{n+1} }{e_{n}}=1-\frac{g(x_n)}{kg(x_n)+e_{n}g'(x_n)}$$
Let's take the limit on both sides of the equation.
$$\lim_{n \to \infty}\frac{e_{n+1} }{e_{n}}=1 -\frac{1}{k}$$
As we know, the convergence rate of Newton's method is quadratic, and in this case, the convergence rate is linear.
So when used Newton's method, we find the convergence rate is linear, we can know that the root is a multiple zero of the function $f$.
(This method is inspired by AI.)\cite{openai2024chatgpt}

\subsubsection*{question2}
similarly, we assume that $f(x)= (x - \alpha  )^{k}g(x)$ and $g(\alpha )\neq 0 $  
$$x_{n+1} - \alpha  = x_n -\alpha - k \frac{(x_{n}-\alpha)g(x_n)}{kg(x_n)+(x_n-\alpha)g'(x_n)}$$
$$x_{n+1} - \alpha = (x_n -\alpha)[1 - k \frac{g(x_n)}{kg(x_n)+(x_n-\alpha)g'(x_n)}]$$
$$e_{n+1}=e_{n}^{2}\frac{g'(x_n)}{kg(x_n)+e_{n}g'(x_n)}$$
$$\lim_{n \to \infty}|\frac{e_{n+1} }{e_{n}^{2}}|=\lim_{n \to \infty}|\frac{g'(x_n)}{kg(x_n)+e_{n}g'(x_n)}|$$
And we know that $e_n \rightarrow 0$ and $x_{n}\rightarrow \alpha $, so we have:
$$\lim_{n \to \infty}|\frac{e_{n+1} }{e_{n}^{2}}|=|\frac{g'(\alpha )}{kg(\alpha )}|$$
and we know that $g(\alpha )\neq 0$ and $g'(\alpha )$ is a const number, so we have proved that the convergence rate is quadratic. $\Box $


% ===============================================
\section*{ \center{\normalsize {Acknowledgement}} }
\printbibliography


\end{document}